\section{Implementación del Protocolo Clásico de Control (Línea Base RSA/ECDSA)}
\label{sec:implementacion_clasica}
% =============================================================================

Para contrastar experimentalmente en el Capítulo 4 el desempeño del protocolo post-cuántico propuesto bajo condiciones idénticas de hardware y arquitectura de software, se implementó un protocolo clásico de referencia que conserva el mismo flujo de validación, las mismas pantallas y la misma estructura de datos, sustituyendo únicamente las primitivas de retículos por los estándares asimétricos tradicionales:

\begin{itemize}
	\item \textbf{Cifrado asimétrico clásico:} Esquema RSA-OAEP con módulo de 2048 bits, exponente público estándar 65537 y función resumen SHA-256 (112 bits de seguridad clásica, vulnerable al algoritmo cuántico de Shor).
	\item \textbf{Firma digital clásica:} Esquema de curvas elípticas ECDSA sobre la curva P-256 del NIST con SHA-256 (128 bits de seguridad clásica, igualmente vulnerable al algoritmo de Shor).
	\item \textbf{Componente simétrico constante:} Se mantiene idéntico el cifrado autenticado AES-256-GCM y el resumen de 256 bits para aislar como única variable independiente el cambio de las primitivas de clave pública.
\end{itemize}

En la Tabla \ref{tab:correspondencia_modulos_didactico} se presenta la equivalencia funcional uno a uno entre los módulos del sistema post-cuántico y los del sistema clásico de control.

\begingroup
\renewcommand{\arraystretch}{1.2}
\setlength{\tabcolsep}{4pt}
\begin{table}[!htbp]
	\centering
	\footnotesize
	\begin{tabular}{|>{\raggedright\arraybackslash}p{3.2cm}|>{\raggedright\arraybackslash}p{5.0cm}|>{\raggedright\arraybackslash}p{5.4cm}|}
		\hline
		\textbf{Componente Lógico} & \textbf{Implementación Post-Cuántica (PQC)} & \textbf{Implementación Clásica de Control} \\
		\hline
		Encapsulamiento / Cifrado & ML-KEM-768 (Estándar FIPS 203) & RSA-OAEP-2048 (Estándar RFC 8017) \\
		\hline
		Esquema de Firma Digital & ML-DSA-65 (Estándar FIPS 204) & ECDSA Curva P-256 (Estándar FIPS 186-4) \\
		\hline
		Núcleo Criptográfico & Implementación propia FIPS 202, 203 y 204 (C/Python) & Módulo clásico sobre OpenSSL (RSA/ECDSA) \\
		\hline
		Entidades del Protocolo & Votante, AR, SV y AC post-cuánticos & Votante, AR, SV y AC clásicos equivalentes \\
		\hline
		Banco de Evaluación & Corredor experimental (300 corridas, lotes de 10 a 1000 votantes) & Corredor experimental bajo idénticas cargas \\
		\hline
		Tamaño Paquete de Voto & \textbf{16293 bytes}* & \textbf{1196 bytes} \\
		\hline
	\end{tabular}
	\caption{Correspondencia entre los módulos post-cuántico y clásico.}
	\label{tab:correspondencia_modulos_didactico}
	\par\vspace{3pt}{\footnotesize\textit{Fuente: Elaboración propia (*medición en formato JSON con campos en hexadecimal y credencial anidada).}}
\end{table}
\endgroup

Esta simetría arquitectónica garantiza que toda variación en los tiempos de ejecución o en el tamaño de los paquetes medida durante las 300 corridas experimentales obedezca exclusivamente al comportamiento algebraico de ML-KEM-768 y ML-DSA-65 frente a RSA-2048 y ECDSA P-256, cuyos resultados estadísticos se analizan en detalle en el Capítulo 4.
